% Printed 391, beginning at the dated heading; preceding header is in the prior component.
\begin{center}
\textsc{Gauss} to \textsc{Gerling}.\quad Göttingen, 12~September 1838.
\end{center}
\par\noindent\hspace*{\parindent}\mbox{.\ .\ .\ .\ .\ .}\quad As for the difference in longitude with Paris, I have never myself compiled the results that \textsc{Triesnecker}, \textsc{Wurm}, and others have obtained from stellar occultations, etc. You yourself are therefore just as well placed as I am to determine \textls[120]{this} result. On the other hand, from the differences in longitude between Jever and Göttingen and likewise between Bentheim and Göttingen, as they followed from my trigonometric measurements and my son's, respectively, and from those of Jever and Bentheim with Paris as given by \textsc{Krayenhoff}, I derived the difference in longitude between Göttingen and Paris that you will find given in \textsc{Harding}'s ephemerides. I have done no new calculation on this account. A complete adjustment of all the triangles that form a ring around the Oldenburg territory would produce a change for Jever, though only a very small one. Jever could even be left out altogether and Emden and Onstwedde used instead, as \textsc{Hartmann}'s measurements reached those places. In any case the change in the final result would be quite insignificant, and I have thus far thought it all the less worth making a new calculation because \textsc{Krayenhoff}'s measurements are, as you can see from my Supplem. theor. comb., much poorer than their reputation.
\par\noindent\hspace*{\parindent}You could obtain a third determination from Altona, whose difference from Göttingen has been determined trigonometrically (see Breitenunterschied [Art.~19]), just as the difference from Greenwich has been determined chronometrically. But I have made no compilation of the relevant data myself and must therefore leave this to you.\quad\mbox{.\ .\ .\ .\ .\ .}

% Printed 392.
\par\clearpage
\noindent\makebox[0pt][l]{\fontsize{13}{16}\selectfont 392}\hfill{\small\textbf{CORRESPONDENCE WITH GERLING.}}\hfill\null\par\medskip
\begin{center}
\textsc{Gauss} to \textsc{Gerling}.\quad Göttingen, 14~November 1838.
\end{center}
\par\noindent\hspace*{\parindent}\mbox{.\ .\ .\ .\ .\ .}\quad You must kindly excuse my having given no further answer to your penultimate letter. I could think of scarcely anything to add to what I communicated in my last letter [of 12~September].
\par\noindent\hspace*{\parindent}As for facts, I really have nothing to add except to refer you to \textsc{Lindenau-Bohnenberger}'s journal, volume IV, p.~119, where you will find the relative position of the old and new observatories here already given. I might add the following two remarks.
\par\noindent\hspace*{\parindent}1) The difference in longitude is given in \textls[120]{time}.
\par\noindent\hspace*{\parindent}2) The position of the new observatory in the passage cited refers to the centre of the rotunda. If I am not mistaken, I previously communicated to you its position relative to the midpoint of the axis of \textsc{Reichenbach}'s meridian circle. If, however, you wish to calculate precisely the position of the latter point relative to the centre of the old observatory, I transcribe the following from my list of coordinates for that purpose:
\begin{center}
\setlength{\tabcolsep}{4pt}
\begin{tabular}{@{}lrr@{}}
Centre of the dome & \(-3{,}104\) & \(-7{,}324\)\\
Midpoint of the axis of \textsc{Reichenbach}'s meridian circle & \(0\) & \(0\)\\
Centre of the old observatory & \(-193{,}54\) & \(+541{,}9\).
\end{tabular}
\end{center}
\par\noindent\hspace*{\parindent}The unit is the metre. In the first number, \(+\) means south, and in the second, \(+\) means west.
\par\noindent\hspace*{\parindent}If you now ask \textls[120]{my} opinion on the \textls[120]{reduction value} to use in expressing your longitudes relative to Paris, I can only say that I regard it as a matter of very little consequence, provided you also state \textls[120]{which} number you took as your basis. Indeed, whatever you choose, the latitude in this respect, or the extent of the leeway, remains
\par\noindent\hspace*{\parindent}\hangindent=2.8em\hangafter=1 A.\quad \textls[120]{extremely} insignificant for any practical purpose, e.g. drawing maps.
\par\noindent\hspace*{\parindent}\hangindent=2.8em\hangafter=1 B.\quad twenty or fifty times greater than all the \textls[120]{relative} differences among all your determinations.

% Printed 393.
\par\clearpage
\noindent\null\hfill{\small\textbf{ON THE HANOVER TRIANGULATION.}}\hfill\makebox[0pt][r]{\fontsize{13}{16}\selectfont 393}\par\medskip
\par\noindent\hspace*{\parindent}Since the latter really constitute the \textls[120]{essence of your work}, this seems to me sufficient justification for the opinion expressed above.
\par\noindent\hspace*{\parindent}In your penultimate letter you speak of the discrepancies in the difference in longitude between Paris and Göttingen obtained by different routes, which amount to as much as \(\frac{3}{4}\) of a second in time, as very considerable. I would not call them that.
\par\noindent\hspace*{\parindent}Indeed, in this connection you seem to me to have entirely overlooked what I worked out on p.~73 [*)] of my paper on the difference in latitude between Göttingen and Altona. If one found geodetically, by two different routes between two points not excessively far apart, differences in longitude that differed by 11 seconds of arc, that would be an enormous error, showing that exceptionally poor observations had crept in or that the calculations had been poorly done.
\par\noindent\hspace*{\parindent}But Paris and Göttingen are very far apart indeed. I believe the measurements of \textsc{Delambre} and \textsc{Krayenhoff} bear no comparison with ours in accuracy, and they were also computed by quite different methods and with quite different elements. I would therefore not be surprised if the difference in longitude thus found differed by some seconds of arc from another also found geodetically, but by a route through, for example, Strassburg, Mannheim, etc.
\par\noindent\hspace*{\parindent}Yet \textls[120]{that} is not at issue here at all. The issue is the difference between the astronomically determined longitudes, which on closer examination mean something \textls[120]{quite different} from geodetic longitudes. The former refer to the progression of the positions of the vertical lines and of the planes drawn through them parallel to the Earth's axis. The latter refer to distances on the surface of the Earth, which \textls[120]{is no ellipsoid at all}, although an ellipsoid can be considered a kind of approximation to it.
\par\noindent\hspace*{\parindent}Given the innumerable known discrepancies between geodetic and astronomical determinations of latitude, one must expect differences of the same order in longitude everywhere. A difference
\par\vfill
\noindent\rule{.22\linewidth}{.4pt}\par\smallskip
{\footnotesize\noindent\hspace*{\parindent}[*) Pp.~49/50 of this volume.]\par}
\medskip\noindent{\small IX.}\hfill 50\hspace*{1.5em}\par

% Printed 394.
\clearpage
\noindent\makebox[0pt][l]{\fontsize{13}{16}\selectfont 394}\hfill{\footnotesize\textbf{CORRESPONDENCE WITH GERLING, ON THE HANOVER TRIANGULATION.}}\hfill\null\par\medskip
\noindent of \(11^{\prime\prime}\), even between two places \textls[120]{much} nearer to one another than Göttingen and Paris, is essentially nothing. There would be absolutely nothing remarkable if even one twice as great were found between Göttingen and Marburg.
\par\noindent\hspace*{\parindent}Do not overlook, then, that the discrepancies you found might very well be entirely correct even if \textls[120]{all} the observations were perfect and absolutely free from error (which they are not). There is \textls[120]{no reason at all} to expect that
\par\noindent\hspace*{\parindent}1) the difference between Göttingen and Paris computed from geodetic measurements on the assumption of particular dimensions of the ellipsoid,
\par\noindent\hspace*{\parindent}2) the one found from astronomical observations,
\par\noindent\hspace*{\parindent}3) the one found by a combination of methods, namely geodetically between Göttingen and Altona and astronomically between Altona-Greenwich-Paris,
\par\noindent should agree with one another, since these three numbers really mean quite different things.
\par\noindent\hspace*{\parindent}If in a future publication of the results of my measurements I should give some of them in the form of differences in latitude and longitude, I will certainly not express the latter with respect to Paris, which does not concern me at all. But if for any reason you think it advisable to communicate yours in such a form, I repeat that I consider it quite immaterial which reduction value you use, so long as the reader is left in no doubt that the fractions of a second appearing in the values can make no claim to exact significance in themselves. They serve only as a means of recovering any desired difference between any two of your numbers with such precision that none of the precision of your measurements themselves is lost. Without wishing to anticipate your decision, I would, on the assumption just stated, take as the basis the number that has come out, taliter qualiter, of the existing geodetic connection between Göttingen and Paris, since the astronomical longitudes, which mean something quite different, are really entirely out of the question here.\quad\mbox{.\ .\ .\ .\ .\ .}
\begin{center}\rule{.12\linewidth}{.4pt}\end{center}

% Printed 395.
\par\clearpage
\noindent\null\hfill{\footnotesize\textbf{REMARKS. CORRESPONDENCE ON THE HANOVER TRIANGULATION.}}\hfill\makebox[0pt][r]{\fontsize{13}{16}\selectfont 395}\par\medskip
\begin{center}\textbf{REMARKS.}\end{center}
\begingroup\small
\par\noindent\hspace*{\parindent}The preceding letters have been printed from the originals.
\par\noindent\hspace*{\parindent}Attached to the letter to \textsc{Bessel} of 26~December 1821, pp.~349/351, was a sketch showing the 5 triangles from the side Hohehagen-Inselsberg to the side Hils-Lichtenberg (see the figure on p.~299) and also the sight lines from Hils towards Deister and Brelingerberg. Brelingerberg, north of Hanover, and Wohlenberg, northwest of Braunschweig, were considered as triangulation stations in the reconnaissance of 1821. At the start of work the following year, however, the former was found unsuitable, and the latter was left out as superfluous. The letter to \textsc{Bessel} of 15~November 1822, pp.~351/358, probably also had a map.
\par\noindent\hspace*{\parindent}Regarding the letters to \textsc{Olbers} of 4 and 8~July 1824, pp.~369/372, it should also be noted that the point Wentel mentioned in them was not included in the system.
\par\noindent\hspace*{\parindent}As \textsc{Gauss}'s report on the work in 1824 to the Hanoverian ministry puts it, »a very good system of triangles had been formed as far as Bremen more quickly than the great difficulties had given reason to hope.« »But the prospects for extending the triangles further now became all the gloomier.\ \mbox{.\ .\ .\ .}\quad After all attempts to connect a new triangulation station to Brüttendorf failed, one expedient remained, namely to include Steinberg in the system of triangles.« But no new triangle could be joined to the side Bremen-Steinberg either. »Without a fortunate circumstance that had just arisen, the prospects for continuing the measurements would have been very poor, since all possibilities now seemed exhausted.« The fortunate circumstance mentioned was \textsc{Gauss}'s discovery from the Ansgariusthurm in Bremen that the tip of the tower at Zeven was still just visible from there. The visibility of Bremen from Zeven had previously gone unnoticed because of peat smoke. Zeven could also be connected not only with Bremen but with the preceding points, Steinberg, Wilsede, and Litberg. The chain of triangles could then be extended from Bremen-Zeven. (See the extract from \textsc{Gauss}'s report on the work in 1825 given later, as well as the letter to \textsc{Olbers} from July 1825, which also follows.)
\par\noindent\hspace*{\parindent}The numerical signals transmitted by heliotrope mentioned in the letter to \textsc{Gerling} of 11~August 1823, pp.~385/386, had prearranged meanings, some of which varied from station to station (according to some brief notes in the observation and calculation books for the arc measurement). (See also the beginning of the letter to \textsc{Gerling} of 1~September 1823, p.~387.) For example, according to one record, 1 meant: things are going well as they are, 2: dim the light, 3: stop for today, 1\,.\,1: depart for the next point.
\par\noindent\hspace*{\parindent}The following letters to \textsc{Schumacher} also concern the arc measurement work in the years 1821/1825: 6 and 30~May, 11~July, 29~September, and 24~October 1821; 10~May, 10~June, 6 and 30~August, 6, 18, 24, and 29~September, 8~October, and 10~November 1822; 8 and 18~June, 23~July, early August, 21~August, 18~September, and 23~October 1823; 24~June, 1~July, 27~September, 17~October, 28~November 1824, and January 1825; 29~April, 14, 20, and 26~June, 11 and 15~July, and 14~August 1825. (See Briefwechsel \textsc{Gauss-Schumacher}, first volume, pp.~229/245, pp.~265/291, pp.~311/325, pp.~334/336, pp.~397/413, pp.~426/430. Second volume, pp.~1/3, pp.~14/31.)
\par\noindent\hspace*{\parindent}The often-mentioned triangulation of the Electorate of Hanover by French engineer-geographers under Lieutenant Colonel \textsc{Epailly} took place in 1804 and 1805 during the French occupation. It was intended to provide the data for the cartographic survey. In early 1821 \textsc{Gauss} received the table of \textsc{Epailly}'s triangles, together with a report by \textsc{Epailly} on his work, from the French Ministry of War through \textsc{Laplace}'s intercession.
\par\noindent\hfill\textsc{Krüger}.\par
\endgroup
\begin{center}\rule{.12\linewidth}{.4pt}\end{center}
\vfill\noindent\hfill 50\textsuperscript{*}\hspace*{1.5em}\par

% Printed 396.
\clearpage
\noindent\makebox[0pt][l]{\fontsize{13}{16}\selectfont 396}\hfill{\small\textbf{ON THE HANOVER TRIANGULATION.}}\hfill\null\par\medskip
\begin{center}
{\large\textbf{PUBLICATIONS}}\par\bigskip
{\large\textbf{ON THE HANOVER TRIANGULATION.}}\par\medskip
[II.]\par\bigskip
\rule{.74\linewidth}{.4pt}\\[-.5ex]
{\small Astronomische Nachrichten, volume I, no.~7, February 1822, pp.~105---106.}\\[-.5ex]
\rule{.74\linewidth}{.4pt}
\end{center}
\par\noindent\hspace*{\parindent}\{It is well known that the Hanoverian government is continuing the arc measurement begun by His Majesty the King of Denmark, that eminent patron of the sciences, which extends in latitude from the northernmost point of Jutland to the southernmost boundary of Lauenburg, and has entrusted the work to Court Councillor \textsc{Gauss} in Göttingen. Our common triangles are Hamburg-Hohenhorn-Lüneburg and Hohenhorn-Lauenburg-Lüneburg. In these triangles the angle at Hamburg was observed by me alone, and the angles at Lüneburg jointly by Court Councillor \textsc{Gauss} and me (the difference between our measurements was under \(0^{\prime\prime}{,}3\)). The angles at Hohenhorn and Lauenburg were measured at the same time (1818) by Captain v.~\textsc{Caroc}. Last year the Court Councillor began his triangles at the southern boundary, and the extract from his letter that I give here contains a brief account of this.
\par\noindent\hfill\textit{S}[\textit{chumacher}].\}\par
\begin{center}Extract from a letter by Court Councillor and Knight \textsc{Gauss}.\end{center}
\par\noindent\hspace*{\parindent}I have the pleasure of showing you the state of my triangulation on the small map enclosed. The heavy lines indicate directions based on measurements already made, and the dotted lines indicate planned directions. I unfortunately have reason to fear that Brelingerberg is visible from neither Wohlenberg nor Lichtenberg (I myself have been to none of the 3 places). In general, forming large triangles in the Lüneburg Heath will be very difficult.

% Printed 397.
\par\clearpage
\noindent\null\hfill{\small\textbf{ON THE HANOVER TRIANGULATION.}}\hfill\makebox[0pt][r]{\fontsize{13}{16}\selectfont 397}\par\medskip
\par\noindent\hspace*{\parindent}Brelingerberg, Deister, Lichtenberg, and Inselsberg have been made visible by heliotrope light, and at the Brocken station so have Hohehagen and Hils, since the signals erected there were only rarely observable, or rather only rarely even possible to \textls[120]{see} (the latter for only a few minutes). With the heliotrope all difficulty disappears.
\par\noindent\hspace*{\parindent}I am now having two others made (of the new design), one of which will soon be finished. I believe I already told you in my last letter that the telegraphic signals I gave from the Brocken using the sextant heliotrope were seen on Hohehagen (distance \(70000\) metres \(=9\frac{1}{3}\) geographical miles) \textls[120]{with the naked eye}. My eldest son Joseph can already set up and aim the earlier heliotrope quite well. With the new one it will probably be a little easier still.
\par\noindent\hspace*{\parindent}The direction from Hils to Hanover has likewise been measured with the utmost precision. Hanover, however, will probably not become a principal triangulation station, since there is only a very limited view to the N.E. from there and Wohlenberg in particular cannot be seen. Moreover, major alterations would first have to be made to the tower there if a theodolite were to be set up on it. Deister will still command an extensive view in all directions, and perhaps even Falkenberg will still be visible there.
\begin{center}
\rule{.74\linewidth}{.4pt}\\[-.5ex]
{\small Astronomische Nachrichten, volume I, no.~24, December 1822, pp.~441---444.}\\[-.5ex]
\rule{.74\linewidth}{.4pt}\par\medskip
Extract from a letter from Court Councillor \textsc{Gauss} to the editor.\\
Göttingen 1822, Nov.~10.
\end{center}
\par\noindent\hspace*{\parindent}Since in the \(7^{\text{th}}\) issue of Astronomische Nachrichten you reported on the state of my triangulation at the close of 1821, I am taking care to send you a brief report on the current state of the operations.
\par\noindent\hspace*{\parindent}Last year the five stations Göttingen, the meridian mark, Hohehagen, Hils, and Brocken were completed, and four points marked out for the further continua-

% Printed 398.
\par\clearpage
\noindent\makebox[0pt][l]{\fontsize{13}{16}\selectfont 398}\hfill{\small\textbf{ON THE HANOVER TRIANGULATION.}}\hfill\null\par\medskip
\noindent tion of operations, namely Lichtenberg, Deister, Wohlenberg, and Brelingerberg. I began this year's work with a reconnaissance trip in the Lüneburg Heath. I considered it all the more necessary because I already knew from the reports of Colonel \textsc{Epailly} the great difficulties of forming a network of triangles in this flat country, which has no appreciable elevations and is covered throughout with woodland in a checkerboard pattern. In 1804 and 1805 he had found these difficulties insurmountable, so he made the connection between Hamburg and southern Hanover by means of a series of triangles along the Weser to its mouth and then back up the Elbe.
\par\noindent\hspace*{\parindent}I found Brelingerberg, which had been sighted from Hils in 1821, \textls[120]{unusable}, as it could not be connected to Lichtenberg and Wohlenberg, but also \textls[120]{superfluous}, like Wohlenberg, because both the site near Garssen and Falkenberg could be connected directly to Lichtenberg. I pass over the great difficulties I had to contend with in extending the triangles farther from Garssen and Falkenberg. These difficulties have now been overcome, and the structure of the network affords many checks, which have turned out to my greatest satisfaction. I will mention only that all my triangulation stations lie at ground level. A stone pedestal of masonry some \(3\frac{1}{2}\)---4 feet high serves to set up the heliotrope and the theodolite. Several lines, in particular those from Falkenberg to Wilsede, from Hauselberg to Breithorn, from Breithorn to Scharnhorst, and from Scharnhorst [to Garssen], required substantial clearings through the woods. Determining the direction of these clearings precisely in advance required elaborate preparations.
\par\noindent\hspace*{\parindent}In the course of the summer I fully completed the stations Lichtenberg, Deister, Garssen, Falkenberg, Hauselberg, Breithorn, Wulfsode, Wilsede, and Scharnhorst, and also made preliminary measurements of the relevant angles on Timpenberg. Thus Hamburg is already provisionally connected, and so is Lüneburg, since you have made a preliminary measurement of the angle between Wilsede and Lüneburg from the Michaelisthurm in Hamburg. Here are some preliminary results, whose absolute scale is provisionally based on von \textsc{Zach}'s baseline near

% Printed 399.
\par\clearpage
\noindent\null\hfill{\small\textbf{ON THE HANOVER TRIANGULATION.}}\hfill\makebox[0pt][r]{\fontsize{13}{16}\selectfont 399}\par\medskip
\noindent Gotha, to which I connected my survey by means of the side from Inselsberg to Brocken.
\begin{center}
\setlength{\tabcolsep}{3pt}
\begin{tabular}{@{}lrrr@{\qquad}rrrl@{}}
 & \multicolumn{3}{c}{\kern-10pt Latitude\kern10pt} & \multicolumn{4}{c}{\kern-5pt\shortstack{Longitude from\\Göttingen}\kern5pt}\\[-1.5ex]
 & \multicolumn{3}{c}{\kern-10pt\(\overbrace{\hphantom{53\quad33\quad31{,}4}}\)\kern10pt} & \multicolumn{4}{c}{\kern-5pt\(\overbrace{\hphantom{0\quad27\quad29{,}5\quad\text{E.}}}\)\kern5pt}\\[-1.2ex]
Hamburg, Michaelisthurm & \(53^\circ\) & \(33'\) & \(1^{\prime\prime}{,}8\) & \(0^\circ\) & \(2'\) & \(3^{\prime\prime}{,}0\) & E.\\
Lüneburg, Michaelisthurm & 53 & 15 & \(5{,}5\) & 0 & 27 & \(29{,}5\) & \multicolumn{1}{c}{»}\\
Celle, southern castle tower & 52 & 37 & \(31{,}4\) & 0 & 8 & \(4{,}9\) & \multicolumn{1}{c}{»}\\
Göttingen observatory, site & & & & & & &\\
\quad of \textsc{Reichenbach}'s & & & & & & &\\
\quad meridian circle & 51 & 31 & \(48{,}7\) & & 0 & &
\end{tabular}
\end{center}
\par\noindent\hspace*{\parindent}The orientation of my system of triangles is taken from my meridian mark. Transferred to Hamburg, it differs by only \(1^{\prime\prime}{,}4\) from the azimuths you communicated to me. To determine the absolute lengths of the lines more precisely, I am only waiting for you to send me the length of your baseline and the triangles that connect it to your principal points. It appears that the one triangle side of mine, Breithorn-Scharnhorst, could be measured directly without insurmountable difficulties.
\par\noindent\hspace*{\parindent}To connect my triangles with yours in a convenient way, I had wished and hoped to link Timpenberg directly to Lüneburg. An attempt was made to cut a clearing, but after it had been carried a considerable distance, the intervening terrain proved too high, and this direct connection therefore had to be abandoned. However, one more point has been established from Wilsede that can be connected directly to Hamburg, Lüneburg, and Lauenburg and most likely to Timpenberg by means of a clearing. The rest must be left to next year's work. No one who has not worked in similar conditions has any idea of the great difficulties of operating in such wooded, flat terrain.
\par\noindent\hspace*{\parindent}The accompanying map, drawn at a scale of \(\frac{1}{2000000}\), will give you a clearer idea of what has been accomplished. Only after the other work was completed did it become apparent that the point Scharnhorst could, by means of two clearings that would not be very difficult to cut, be connected directly

% Printed 400.
\par\clearpage
\noindent\makebox[0pt][l]{\fontsize{13}{16}\selectfont 400}\hfill{\small\textbf{ON THE HANOVER TRIANGULATION.}}\hfill\null\par\medskip
\noindent to Lichtenberg and Deister. Had it been possible to find this site earlier and determine its suitability and position, Garssen could have been omitted altogether. Perhaps I shall still measure the angles of the triangle Scharnhorst-Deister-Lichtenberg next year.
\par\noindent\hspace*{\parindent}This year, besides the heliotrope used in 1821, I have had two others of the new design in operation, as well as another heliotrope apparatus, which I always carried with me to give my assistants orders by telegraph. I need scarcely tell you that the account of the design of heliotropes given by Mr.~\textsc{Schubach} in the Astron. Jahrbuch for 1825 rests entirely on an error and has nothing in common with \textls[120]{my} heliotropes. Mr.~\textsc{Rumpf} has already made seven heliotropes, of which two are intended for the Prussian triangulation and two for the Hessian triangulation. Drawings of both designs are at your disposal on request.
\begin{center}\rule{.12\linewidth}{.4pt}\end{center}
\begin{center}\textbf{REMARKS.}\end{center}
\begingroup\small
\par\noindent\hspace*{\parindent}The heavy lines on the triangulation map accompanying the first article in volume I, no.~7, of Astronomische Nachrichten show in the south the triangles from the side Hohehagen-Inselsberg to the side Hils-Lichtenberg (see the figure on p.~299) and the sight lines from Hils to Deister and Brelingerberg, and in the north the two triangles Lüneburg-Hamburg-Hohenhorn and Lüneburg-Hohenhorn-Lauenburg. Dotted lines mark the triangle Lichtenberg-Deister-Wohlenberg, the directions between the points Deister, Brelingerberg, Falkenberg, and Wilsede, and the directions Wilsede-Lüneburg and Wilsede-Hamburg. The line Hils-Northern meridian mark is missing from this map, presumably by oversight.
\par\noindent\hspace*{\parindent}The triangulation map attached to the second article shows the triangles of the Hanoverian arc measurement proper, that is, those shown in the figure on p.~299 from Inselsberg to Hamburg, extending to the sides Falkenberg-Wilsede and Wilsede-Hamburg. The point Nindorf and its connections are still missing from it, whereas it includes the connecting triangles between the points Lüneburg, Wilsede, Hamburg, and Hohenhorn, as well as the two Danish points Syk and Lauenburg.
\par\noindent\hspace*{\parindent}Besides these two communications on the Hanover triangulation, \textsc{Gauss} published another 2 papers from 1821 on the »Heliotrop« (which follow later), and an article from 1823, »Beobachtete und berechnete Triangulirung im Hannoverschen, Braunschweigschen und Lüneburgschen« (also given later).
\par\noindent\hfill\textsc{Krüger}.\par
\endgroup
\begin{center}\rule{.16\linewidth}{.4pt}\end{center}
% End of printed 400. Printed 401 opens III. NACHLASS ZUR HANNOVERSCHEN TRIANGULATION.
